\documentclass[a4paper,12pt,twoside]{memoir}

% Castellano
\usepackage[spanish,es-tabla]{babel}
\selectlanguage{spanish}
\usepackage[utf8]{inputenc}
\usepackage[T1]{fontenc}
\usepackage{lmodern} % Scalable font
\usepackage{microtype}
\usepackage{placeins}

\RequirePackage{booktabs}
\RequirePackage[table]{xcolor}
\RequirePackage{xtab}
\RequirePackage{multirow}

% Links
\PassOptionsToPackage{hyphens}{url}\usepackage[colorlinks]{hyperref}
\hypersetup{
	allcolors = {red}
}

% Ecuaciones
\usepackage{amsmath}

% Rutas de fichero / paquete
\newcommand{\ruta}[1]{{\sffamily #1}}

% Párrafos
\nonzeroparskip

% Huérfanas y viudas
\widowpenalty100000
\clubpenalty100000

% Imágenes

% Comando para insertar una imagen en un lugar concreto.
% Los parámetros son:
% 1 --> Ruta absoluta/relativa de la figura
% 2 --> Texto a pie de figura
% 3 --> Tamaño en tanto por uno relativo al ancho de página
\usepackage{graphicx}
\newcommand{\imagen}[3]{
	\begin{figure}[!h]
		\centering
		\includegraphics[width=#3\textwidth]{#1}
		\caption{#2}\label{fig:#1}
	\end{figure}
	\FloatBarrier
}

% Comando para insertar una imagen sin posición.
% Los parámetros son:
% 1 --> Ruta absoluta/relativa de la figura
% 2 --> Texto a pie de figura
% 3 --> Tamaño en tanto por uno relativo al ancho de página
\newcommand{\imagenflotante}[3]{
	\begin{figure}
		\centering
		\includegraphics[width=#3\textwidth]{#1}
		\caption{#2}\label{fig:#1}
	\end{figure}
}

% El comando \figura nos permite insertar figuras comodamente, y utilizando
% siempre el mismo formato. Los parametros son:
% 1 --> Porcentaje del ancho de página que ocupará la figura (de 0 a 1)
% 2 --> Fichero de la imagen
% 3 --> Texto a pie de imagen
% 4 --> Etiqueta (label) para referencias
% 5 --> Opciones que queramos pasarle al \includegraphics
% 6 --> Opciones de posicionamiento a pasarle a \begin{figure}
\newcommand{\figuraConPosicion}[6]{%
  \setlength{\anchoFloat}{#1\textwidth}%
  \addtolength{\anchoFloat}{-4\fboxsep}%
  \setlength{\anchoFigura}{\anchoFloat}%
  \begin{figure}[#6]
    \begin{center}%
      \Ovalbox{%
        \begin{minipage}{\anchoFloat}%
          \begin{center}%
            \includegraphics[width=\anchoFigura,#5]{#2}%
            \caption{#3}%
            \label{#4}%
          \end{center}%
        \end{minipage}
      }%
    \end{center}%
  \end{figure}%
}

%
% Comando para incluir imágenes en formato apaisado (sin marco).
\newcommand{\figuraApaisadaSinMarco}[5]{%
  \begin{figure}%
    \begin{center}%
    \includegraphics[angle=90,height=#1\textheight,#5]{#2}%
    \caption{#3}%
    \label{#4}%
    \end{center}%
  \end{figure}%
}
% Para las tablas
\newcommand{\otoprule}{\midrule [\heavyrulewidth]}
%
% Nuevo comando para tablas pequeñas (menos de una página).
\newcommand{\tablaSmall}[5]{%
 \begin{table}
  \begin{center}
   \rowcolors {2}{gray!35}{}
   \begin{tabular}{#2}
    \toprule
    #4
    \otoprule
    #5
    \bottomrule
   \end{tabular}
   \caption{#1}
   \label{tabla:#3}
  \end{center}
 \end{table}
}

%
% Nuevo comando para tablas pequeñas (menos de una página).
\newcommand{\tablaSmallSinColores}[5]{%
 \begin{table}[H]
  \begin{center}
   \begin{tabular}{#2}
    \toprule
    #4
    \otoprule
    #5
    \bottomrule
   \end{tabular}
   \caption{#1}
   \label{tabla:#3}
  \end{center}
 \end{table}
}

\newcommand{\tablaApaisadaSmall}[5]{%
\begin{landscape}
  \begin{table}
   \begin{center}
    \rowcolors {2}{gray!35}{}
    \begin{tabular}{#2}
     \toprule
     #4
     \otoprule
     #5
     \bottomrule
    \end{tabular}
    \caption{#1}
    \label{tabla:#3}
   \end{center}
  \end{table}
\end{landscape}
}

%
% Nuevo comando para tablas grandes con cabecera y filas alternas coloreadas en gris.
\newcommand{\tabla}[6]{%
  \begin{center}
    \tablefirsthead{
      \toprule
      #5
      \otoprule
    }
    \tablehead{
      \multicolumn{#3}{l}{\small\sl continúa desde la página anterior}\\
      \toprule
      #5
      \otoprule
    }
    \tabletail{
      \hline
      \multicolumn{#3}{r}{\small\sl continúa en la página siguiente}\\
    }
    \tablelasttail{
      \hline
    }
    \bottomcaption{#1}
    \rowcolors {2}{gray!35}{}
    \begin{xtabular}{#2}
      #6
      \bottomrule
    \end{xtabular}
    \label{tabla:#4}
  \end{center}
}

%
% Nuevo comando para tablas grandes con cabecera.
\newcommand{\tablaSinColores}[6]{%
  \begin{center}
    \tablefirsthead{
      \toprule
      #5
      \otoprule
    }
    \tablehead{
      \multicolumn{#3}{l}{\small\sl continúa desde la página anterior}\\
      \toprule
      #5
      \otoprule
    }
    \tabletail{
      \hline
      \multicolumn{#3}{r}{\small\sl continúa en la página siguiente}\\
    }
    \tablelasttail{
      \hline
    }
    \bottomcaption{#1}
    \begin{xtabular}{#2}
      #6
      \bottomrule
    \end{xtabular}
    \label{tabla:#4}
  \end{center}
}

%
% Nuevo comando para tablas grandes sin cabecera.
\newcommand{\tablaSinCabecera}[5]{%
  \begin{center}
    \tablefirsthead{
      \toprule
    }
    \tablehead{
      \multicolumn{#3}{l}{\small\sl continúa desde la página anterior}\\
      \hline
    }
    \tabletail{
      \hline
      \multicolumn{#3}{r}{\small\sl continúa en la página siguiente}\\
    }
    \tablelasttail{
      \hline
    }
    \bottomcaption{#1}
  \begin{xtabular}{#2}
    #5
   \bottomrule
  \end{xtabular}
  \label{tabla:#4}
  \end{center}
}



\definecolor{cgoLight}{HTML}{EEEEEE}
\definecolor{cgoExtralight}{HTML}{FFFFFF}

%
% Nuevo comando para tablas grandes sin cabecera.
\newcommand{\tablaSinCabeceraConBandas}[5]{%
  \begin{center}
    \tablefirsthead{
      \toprule
    }
    \tablehead{
      \multicolumn{#3}{l}{\small\sl continúa desde la página anterior}\\
      \hline
    }
    \tabletail{
      \hline
      \multicolumn{#3}{r}{\small\sl continúa en la página siguiente}\\
    }
    \tablelasttail{
      \hline
    }
    \bottomcaption{#1}
    \rowcolors[]{1}{cgoExtralight}{cgoLight}

  \begin{xtabular}{#2}
    #5
   \bottomrule
  \end{xtabular}
  \label{tabla:#4}
  \end{center}
}



\graphicspath{ {./img/} }

% Capítulos
\chapterstyle{bianchi}
\newcommand{\capitulo}[2]{
	\setcounter{chapter}{#1}
	\setcounter{section}{0}
	\setcounter{figure}{0}
	\setcounter{table}{0}
	\chapter*{\thechapter.\enskip #2}
	\addcontentsline{toc}{chapter}{\thechapter.\enskip #2}
	\markboth{#2}{#2}
}

% Apéndices
\renewcommand{\appendixname}{Apéndice}
\renewcommand*\cftappendixname{\appendixname}

\newcommand{\apendice}[1]{
	%\renewcommand{\thechapter}{A}
	\chapter{#1}
}

\renewcommand*\cftappendixname{\appendixname\ }

% Formato de portada
\makeatletter
\usepackage{xcolor}
\newcommand{\tutor}[1]{\def\@tutor{#1}}
\newcommand{\course}[1]{\def\@course{#1}}
\definecolor{cpardoBox}{HTML}{E6E6FF}
\def\maketitle{
  \null
  \thispagestyle{empty}
  % Cabecera ----------------
\noindent\includegraphics[width=\textwidth]{cabecera}\vspace{1cm}%
  \vfill
  % Título proyecto y escudo informática ----------------
  \colorbox{cpardoBox}{%
    \begin{minipage}{.8\textwidth}
      \vspace{.5cm}\Large
      \begin{center}
      \textbf{TFG del Grado en Ingeniería Informática}\vspace{.6cm}\\
      \textbf{\LARGE\@title{}}
      \end{center}
      \vspace{.2cm}
    \end{minipage}

  }%
  \hfill\begin{minipage}{.20\textwidth}
    \includegraphics[width=\textwidth]{escudoInfor}
  \end{minipage}
  \vfill
  % Datos de alumno, curso y tutores ------------------
  \begin{center}%
  {%
    \noindent\LARGE
    Presentado por \@author{}\\ 
    en Universidad de Burgos --- \@date{}\\
    Tutor: \@tutor{}\\
  }%
  \end{center}%
  \null
  \cleardoublepage
  }
\makeatother

\newcommand{\nombre}{Nombre del alumno} %%% cambio de comando

% Datos de portada
\title{título del TFG}
\author{\nombre}
\tutor{nombre tutor}
\date{\today}

\begin{document}

\maketitle


\newpage\null\thispagestyle{empty}\newpage


%%%%%%%%%%%%%%%%%%%%%%%%%%%%%%%%%%%%%%%%%%%%%%%%%%%%%%%%%%%%%%%%%%%%%%%%%%%%%%%%%%%%%%%%
\thispagestyle{empty}


\noindent\includegraphics[width=\textwidth]{cabecera}\vspace{1cm}

\noindent D. nombre tutor, profesor del departamento de nombre departamento, área de nombre área.

\noindent Expone:

\noindent Que el alumno D. \nombre, con DNI dni, ha realizado el Trabajo final de Grado en Ingeniería Informática titulado título de TFG. 

\noindent Y que dicho trabajo ha sido realizado por el alumno bajo la dirección del que suscribe, en virtud de lo cual se autoriza su presentación y defensa.

\begin{center} %\large
En Burgos, {\large \today}
\end{center}

\vfill\vfill\vfill

% Author and supervisor
\begin{minipage}{0.45\textwidth}
\begin{flushleft} %\large
Vº. Bº. del Tutor:\\[2cm]
D. nombre tutor
\end{flushleft}
\end{minipage}
\hfill
\begin{minipage}{0.45\textwidth}
\begin{flushleft} %\large
Vº. Bº. del co-tutor:\\[2cm]
D. nombre co-tutor
\end{flushleft}
\end{minipage}
\hfill

\vfill

% para casos con solo un tutor comentar lo anterior
% y descomentar lo siguiente
%Vº. Bº. del Tutor:\\[2cm]
%D. nombre tutor


\newpage\null\thispagestyle{empty}\newpage




\frontmatter

% Abstract en castellano
\renewcommand*\abstractname{Resumen}
\begin{abstract}
En este primer apartado se hace una \textbf{breve} presentación del tema que se aborda en el proyecto.
\end{abstract}

\renewcommand*\abstractname{Descriptores}
\begin{abstract}
Palabras separadas por comas que identifiquen el contenido del proyecto Ej: servidor web, buscador de vuelos, android \ldots
\end{abstract}

\clearpage

% Abstract en inglés
\renewcommand*\abstractname{Abstract}
\begin{abstract}
A \textbf{brief} presentation of the topic addressed in the project.
\end{abstract}

\renewcommand*\abstractname{Keywords}
\begin{abstract}
keywords separated by commas.
\end{abstract}

\clearpage

% Indices
\tableofcontents

\clearpage

\listoffigures

\clearpage

\listoftables
\clearpage

\mainmatter
\capitulo{1}{Introducción}

Descripción del contenido del trabajo y del estructura de la memoria y del resto de materiales entregados.

\capitulo{2}{Objetivos del proyecto}

Este apartado explica de forma precisa y concisa cuales son los objetivos que se persiguen con la realización del proyecto. Se puede distinguir entre los objetivos marcados por los requisitos del software a construir y los objetivos de carácter técnico que plantea a la hora de llevar a la práctica el proyecto.

\capitulo{3}{Conceptos teóricos}

En este apartado se definen algunos términos y conceptos relevantes sobre el proyecto.

\subsection{Soft Skills}
Las habilidades blandas ~\cite{wiki:SoftSkills}, competencias blandas o habilidades suaves  son una combinación de habilidades sociales, habilidades de comunicación, rasgos de la personalidad, actitudes, atributos profesionales, inteligencia social e inteligencia emocional, que facultan a las personas para moverse por su entorno, trabajar bien con otros, realizar un buen desempeño y, complementándose con las habilidades duras, conseguir sus objetivos.
Las habilidades blandas son un cúmulo de rasgos productivos de la personalidad que caracterizan las relaciones de una persona en un medio. Estas habilidades pueden incluir autoestima, comunicación, elocuencia, hábitos personales, empatía, gestión del tiempo, trabajo en equipo y liderazgo. Una definición basada en la revisión de artículos al respecto considera "habilidades blandas" un término paraguas para habilidades con tres elementos funcionales clave: interpersonales, sociales y profesionales.
Las habilidades blandas se pueden categorizar en diferentes como estas:
\begin{enumerate}
    \item \textbf{Comunicación}: capacidad de expresarse claramente hablando y escribiendo, y de hacer presentaciones en público.
    \item \textbf{Cortesía}: buenos modales, etiqueta, respeto, decir «por favor» y «gracias».
    \item \textbf{Flexibilidad}: adaptabilidad, disposición al cambio y a la formación continua, aceptación de lo nuevo, se le puede enseñar.
    \item \textbf{Integridad}: sinceridad, moralidad, valores personales, honradez.
    \item \textbf{Habilidades interpersonales}: agradable, con sentido del humor, amistoso, hospitalario, empático, con autocontrol, paciencia, sociabilidad, calidez, habilidades sociales.
    \item \textbf{Actitud positiva}: optimista, entusiasta, ánima a sus compañeros, feliz, seguro.
    \item \textbf{Profesionalidad}: bien vestido, buen aspecto, compostura.
    \item \textbf{Responsabilidad}: responsable, fiable, termina el trabajo, quiere hacerlo bien, es consciente, tiene recursos, autodisciplina y sentido común.
    \item \textbf{Trabajo en equipo}: coopera, se lleva bien con otros, es agradable, anima, ayuda a quien lo necesita.
    \item\textbf{Actitud hacia el trabajo}: trabaja duro, está disponible, es leal, puntual, tiene iniciativa, dedicación, motivación y falta poco al trabajo.
    \item \textbf{Creatividad} ~\cite{wiki:Creatividad}: capacidad de crear nuevas ideas o conceptos, de nuevas asociaciones entre ideas y conceptos conocidos, que habitualmente producen soluciones originales.
    \item \textbf{Resolución de problemas} ~\cite{wiki:ResProblemas}: la conclusión de un proceso más amplio que tiene como pasos previos la identificación del problema y su modelado.
    \item \textbf{Ser organizado}: llevar a cabo unos pasos de manera seguida non un propósito aprovechando el tiempo lo mas posible.
\end{enumerate}
Estas pueden ser utilizadas para poder hacer un cribado de un gran grupo de personas y así poder realizar una agrupación con diferentes integrantes que tengas diferentes atributos como el de líder, trabajador, creativo, espontaneo, observador, etc. ; o por otra parte hacer una selección para un nuevo proyecto a partir de una semejanza a los conocimientos deseados.
~\cite{softskillsacademy}
\imagen{softSkills}{Tabla de Soft Skills}{1}
\capitulo{4}{Técnicas y herramientas}

Esta parte de la memoria tiene como objetivo presentar las técnicas metodológicas y las herramientas de desarrollo que se han utilizado para llevar a cabo el proyecto. Si se han estudiado diferentes alternativas de metodologías, herramientas, bibliotecas se puede hacer un resumen de los aspectos más destacados de cada alternativa, incluyendo comparativas entre las distintas opciones y una justificación de las elecciones realizadas. 
No se pretende que este apartado se convierta en un capítulo de un libro dedicado a cada una de las alternativas, sino comentar los aspectos más destacados de cada opción, con un repaso somero a los fundamentos esenciales y referencias bibliográficas para que el lector pueda ampliar su conocimiento sobre el tema.



\capitulo{5}{Aspectos relevantes del desarrollo del proyecto}

En este apartado indicare las elecciones que he tomado, el transcurro del proyecto y como se han abordado algún problema relevante.
\section{Elección de las herramientas}
Una vez se empieza a trabajar hay que poner los cimientos, por lo que hay que realizar una selección cuidadosa con los elementos que se va a realizar el proyecto. Por eso con la ayuda de esta \href{ https://stackshare.io/stackups/trending }{ herramienta web} para comparar aspectos de herramientas con el mismo propósito he podido hacer una elección cómoda. 
\subsection{Visual Studio Code}
Para la elección del editor de texto tenemos varias opciones como pueden ser “PyCharm,, Atom, Komodo y VSCodium” todas estas opciones son válidas pero si tenemos que elegir una sería PyCharm la cual no he escogido debido al haber ya trabajado repetidamente con Visual Studio Code y estar más familiarizado, en parte para las necesidades que necesito no hay ninguna extensión que tenga que implementar que no venga ya incluida con Visual Studio Code.
Diferencias entre Visual Studio Code y PyCharm utilizadas para la selección ~\cite{VSC-PyCharm}:
\begin{itemize}
    \item \textbf{Interfaz de usuario:} PyCharm ofrece una interfaz de usuario más completa y con más funciones, que ofrece una experiencia de IDE rica y profesional con herramientas avanzadas de depuración y creación de perfiles. Por otro lado, Visual Studio Code presenta una interfaz de usuario liviana y minimalista, que se centra en la simplicidad y la extensibilidad, con una amplia gama de extensiones disponibles y opciones de personalización.
    \item \textbf{Edición de código:} PyCharm ofrece una potente asistencia de codificación y finalización de código inteligente, brindando sugerencias útiles e importaciones automáticas para módulos de Python. También incluye funciones especializadas para Django, Flask y otros marcos. Si bien Visual Studio Code también brinda finalización de código y sugerencias, es posible que requiera extensiones adicionales para igualar el nivel de inteligencia de código que brinda PyCharm.
    \item \textbf{Herramientas y complementos integrados:} PyCharm incluye un conjunto completo de herramientas integradas para el control de versiones (por ejemplo, Git), marcos de prueba, bases de datos y desarrollo web. Ofrece una integración perfecta con herramientas populares de Python, como entornos virtuales y administradores de paquetes. Visual Studio Code, por otro lado, depende en gran medida del uso de complementos y extensiones para proporcionar funcionalidades similares, lo que permite a los desarrolladores personalizar su entorno con las herramientas que elijan.
    \item \textbf{Rendimiento:} PyCharm está basado en la plataforma IntelliJ, que puede consumir una cantidad significativa de recursos del sistema y generar tiempos de inicio e indexación más lentos en comparación con Visual Studio Code. Visual Studio Code, al ser un IDE liviano, generalmente ofrece un rendimiento más rápido y tiempos de inicio más rápidos.
    \item \textbf{Comunidad y ecosistema:} PyCharm cuenta con una comunidad consolidada y madura con una amplia documentación y foros de soporte. Es ampliamente utilizado en la comunidad de desarrollo de Python, lo que se traduce en una gran cantidad de recursos y tutoriales disponibles. Visual Studio Code, desarrollado por Microsoft, se beneficia de su amplia comunidad y ecosistema, y ofrece una amplia gama de extensiones, temas e integraciones impulsados por la comunidad con otras herramientas de Microsoft.
    \item \textbf{Precios:} PyCharm ofrece ediciones comunitarias tanto gratuitas como de pago. Las versiones de pago ofrecen funciones adicionales y soporte para el desarrollo web y empresarial. Visual Studio Code, por otro lado, es completamente gratuito y de código abierto, lo que lo convierte en una opción atractiva para los desarrolladores que prefieren una solución rentable.
\end{itemize}

\subsection{GitHub}
En cuanto a la selección de la plataforma donde alojar el proyecto GitHub es muy común, tras realizar diferentes proyectos en esta plataforma ya me encuentro familiarizado y no he tenido ninguna intención de realizar el alojamiento en otras plataformas. Por otra parte la integración de la metodología Kanban por parte de Zube resulta más sencillo el manipular las tareas.
Diferencias entre GitHub y GitLab utilizadas para la selección ~\cite{GitHub-GitLab}:
\begin{itemize}
    \item \textbf{Alojamiento e implementación:} GitHub es principalmente una plataforma basada en la nube que ofrece alojamiento para repositorios Git. Proporciona un servicio totalmente administrado, donde los usuarios pueden crear y administrar sus repositorios directamente en los servidores de GitHub. GitLab, por otro lado, ofrece una solución SaaS basada en la nube (GitLab.com) y una solución alojada por los propios usuarios (GitLab Community Edition). Esto significa que los usuarios tienen la flexibilidad de elegir entre alojar sus repositorios en la plataforma en la nube de GitLab o implementar GitLab en su propia infraestructura.
    \item \textbf{Precios y licencias:} GitHub ofrece opciones gratuitas para repositorios públicos y planes pagos para repositorios privados, mientras que GitLab ofrece una Community Edition gratuita y una Enterprise Edition paga.
    \item \textbf{Conjunto de características:} GitHub es conocido por su sólido ecosistema de integración, que incluye una amplia gama de herramientas y servicios de terceros para la integración continua, la implementación y la gestión de proyectos. GitLab, por otro lado, tiene como objetivo proporcionar una plataforma DevOps más completa con características integradas como canalizaciones de CI/CD, integración con Kubernetes y un registro de contenedores Docker integrado.
    \item \textbf{Comunidad y colaboración:} GitHub cuenta con una comunidad de desarrolladores grande y activa, lo que la convierte en una plataforma popular para proyectos de código abierto. Ofrece funciones que facilitan la colaboración y la contribución, como solicitudes de incorporación de cambios, revisiones de código y seguimiento de problemas. GitLab también cuenta con una comunidad en crecimiento y su naturaleza de código abierto fomenta la colaboración. Además, GitLab enfatiza el concepto de "Flujo de GitLab", que promueve un flujo de trabajo más integrado y optimizado para el desarrollo, la colaboración y la implementación. Esto incluye funciones como canales de CI/CD integrados y una interfaz consolidada para administrar el código, los problemas y la gestión de proyectos.
\end{itemize}

\subsection{Flask}
Para la realización de la página web una manera sencilla de empezar en el caso de no haber realizado ninguna anteriormente, es Flask por lo cual entre la propia ayuda que te ofrece la documentación de flask y el gran abanico de librerías que ofrece son suficientes para el proyecto que tenía en mente.
Diferencias entre GitHub y GitLab utilizadas para la selección ~\cite{Flask-Django}:
\begin{itemize}
    \item \textbf{Tipo de framework:} Django es un framework web de alto nivel y de pila completa que sigue el patrón de diseño Modelo-Vista-Plantilla (MVT), y ofrece funciones integradas como autenticación, interfaz de administración y ORM. Por otro lado, Flask es un framework web ligero que se basa en Werkzeug y Jinja2 y ofrece flexibilidad a los desarrolladores para elegir las herramientas y bibliotecas con las que quieren trabajar.
    \item \textbf{Escalabilidad:} Django es más adecuado para aplicaciones más grandes debido a sus características y convenciones integradas que ayudan a gestionar proyectos complejos de manera eficiente. Flask, al ser más minimalista, es más adecuado para proyectos más pequeños o aplicaciones que requieren un enfoque más personalizado.
    \item \textbf{Soporte y ecosistema de la comunidad:} Django tiene una comunidad más grande y madura en comparación con Flask, y ofrece una amplia gama de paquetes de terceros, tutoriales y documentación que pueden ayudar a los desarrolladores a crear e implementar aplicaciones más rápido. Flask, por otro lado, tiene una comunidad más pequeña, pero ofrece un ecosistema más liviano y flexible que permite a los desarrolladores tener más control sobre sus proyectos.
    \item \textbf{Curva de aprendizaje:} Django tiene una curva de aprendizaje más pronunciada en comparación con Flask debido a sus características integrales y convenciones que los desarrolladores deben comprender y seguir. Flask, al ser más simple y minimalista, tiene una barrera de entrada más baja, lo que facilita que los principiantes comiencen a crear aplicaciones web.
    \item \textbf{Enrutamiento de URL:} En Django, el enrutamiento de URL se realiza a través de una configuración de URL predefinida mediante expresiones regulares, lo que lo hace más estructurado y organizado. Flask, por otro lado, utiliza decoradores para definir rutas, lo que proporciona una forma más flexible e intuitiva de enrutar URL en la aplicación.
    \item \textbf{Motor de plantillas:} Django viene con un motor de plantillas integrado que sigue los principios DRY (Don't Repeat Yourself), lo que permite a los desarrolladores reutilizar el código de las plantillas de manera eficiente. Flask, por otro lado, utiliza Jinja2 como su motor de plantillas predeterminado, lo que proporciona una sintaxis más liviana y fácil de usar para el diseñador para crear plantillas.
\end{itemize}

\section{Desarrollo del proyecto}
Para el desarrollo he seguido unas pautas que han sido claves, las cuales consisten en:
\begin{enumerate}
    \item Plantear la idea a fondo con un propósito adicional al deseado para abarcar posibles cambios.
    \item Realizar un primer código con el cual se logre lo deseado y después adaptarlo para obtener la funcionalidad deseada.
    \item Realizar una prueba de todo el programa para depurar y comprobar que no haya ningún error.
\end{enumerate}

\section{Problemas}
Los problemas  siempre aparecen y como no podía ser durante el desarrollo de la aplicación surgieron diferentes errores los cuales se han ido desarrollado. 

Un problema que tuve fue la visualización de caracteres erróneos cuando el texto disponía de alguna tilde o símbolo de interrogación, al estar trabajando con todos los ficheros en “UTF-8” comprobándolo con el propio “Visual Studio Code” o el “notepad++” y de asegurar que el texto del html estuviese codificado también adecuadamente,  no tenía el porqué del error, pero tras darle algunas vueltas se me ocurrió la idea de que no estuviese tratando el fichero correctamente y de esta manera lo solucione al especificar como quería abrir el fichero en cuestión para obtener los datos ya así codificados en “UTF-8”. 

Durante la implementación del usuario invitado para poder realizar la evaluación sin estar registrado, me encontré con que en un principio al solo tener la evaluación de un test al añadir el modelo de los cuestionarios tuve el dilema de cómo poner en común al invitado, por lo que para poder solucionarlo cree una función intermedia en la cual comprobaba si estaba logeado algún usuario en la sesión así pudiendo  utilizar el mismo invitado en común para luego llamar a la función correspondiente dependiendo si se quiere hacer un test o cuestionario. 

A la hora de utilizar Docker para crear un contenedor de la aplicación me surgió un error en el cual  las imágenes correspondientes al test no se muestran dando como fallo que no se encuentran, el problema en cuestión resultaba en que las imágenes no se encontraban en el destino pero viendo los archivos del contenedor las imágenes si se encontraban, lo cual para solucionarlo probé diferentes maneras de llegar a la ruta pero al final el error en si es que los propios archivos tenían la extensión  \textbf{.PNG} mientras utilizaba la extensión \textbf{.png} como ruta para cargar las imágenes, como docker genera una máquina virtual con un sistema operativo de Linux, al tratar los archivos las extensiones son importantes, por lo que para solucionarlo tuve que cambiar la extensión de la ruta que apunta a las imágenes del json a \textbf{.PNG}.
\capitulo{6}{Trabajos relacionados}

Este apartado sería parecido a un estado del arte de una tesis o tesina. En un trabajo final grado no parece obligada su presencia, aunque se puede dejar a juicio del tutor el incluir un pequeño resumen comentado de los trabajos y proyectos ya realizados en el campo del proyecto en curso. 

\capitulo{7}{Conclusiones y Líneas de trabajo futuras}

Para finalizar se expondrán las conclusiones finales obtenidas a la realización del proyecto y por otra parte las líneas de trabajo futuras.

\subsection{Conclusiones}
A medida que se avanzaba en el proyecto los objetivos marcados se realizaban satisfactoriamente en su mayoría aun al no haber podido realizar alguna funcionalidad de la manera deseada.
La navegación entre las pestañas se hace intuitiva y no es pesado lo cual nos proporciona una página web que es accesible y no necesita tener una guía a lado para poder interactuar con ella satisfactoriamente.
Al realizar el proyecto se ha obtenido una perspectiva que no se tenía respecto a cómo desarrollar una página web y diferentes funcionalidades que debería tener de base.
Como resumen la realización de este proyecto pone a prueba la utilización de diferentes aspectos enseñados en la carrera, además que se hace un planteamiento único que resuelve los diferentes problemas que van surgiendo a medida. 
\subsection{Líneas de trabajo futuras} 
Aunque se tenga una funcionalidad correcta y unas funcionalidades satisfechas, siempre es posible realizar mejoras para tener un producto mejor.
\begin{itemize}
    \item Mejora de la lógica de cambio de contraseña aplicando una doble verificación al enviar un mensaje al correo asociado a la cuenta que desea cambiar la contraseña.
    \item Implementación de una opción que te permita modificar los cuestionarios/tests o crear nuevos desde la propia aplicación web.
    \item Implementación de un nuevo rol “editor” que sea el encargado del apartado anterior.
    \item Poder hacer un filtrado sobre los diferentes campos respecto a las tablas de usuarios y los resultados.
    \item Creación de unas pruebas unitarias para comprobar un funcionamiento correcto sin tener que realizarlo a mano.
\end{itemize}



\bibliographystyle{plain}
\bibliography{bibliografia}

\end{document}
